\chapter{Halbwertsbreiten}

Die Halbwertsbreite (FWHM) von Gaußkurven ist gegeben über \todo{source}
\begin{align}
	\text{FWHM} = \sigma \cdot \sqrt{8 \ln(2)} \label{eq:fwhm}
\end{align}

Interessant sind die Halbwertsbreiten in der Energie-Domäne. Da die Gaußanpassungen in Kanälen durchgeführt werden, erfolgt die Umrechnung mit $a$, der Steigung aus der Energiekalibration.

\begin{align}
	\Delta E = \text{FWHM}_E = \text{FWHM}_b * a \label{eq:fwhme}
\end{align}


Die FWHM der angepassten Linien aus Tabelle \todo{refrence} werden nun mit Gl \eqref{eq:fwhm} und \eqref{eq:fwhme} berechnet. Die Ergebnisse sind in Tabelle \todo{refrence} zu finden. Man erkennt sofort, dass die Halbwertsbreite des HPGE Detektors um mehrere Größenordnungen kleiner ist, als die des NaI Detektors.

Anhand der Identifizierten Europiumlinien im HPGE Detektor wird der intrinsische- und elektronische Teil der Halbwertsbreite bestimmt. Für diese gilt der Zusammenhang

\begin{align}
	\Delta E (E_\gamma) &= \sqrt{(\Delta E_d (E_\gamma))^2 + (\Delta E_e)^2}\\
	\Delta E_d (E_\gamma) &= k \cdot \sqrt{E_\gamma}
\end{align}

wobei $E_d$ die intrinsische und $E_e$ die elektronische Halbwertsbreite sind \cite[5-6]{521versuchsanleitung}. 
Zum Überprüfen dieses Zusammenhangs wird 

\begin{align}
	(\Delta E (E_\gamma))^2 &= (\Delta E_d (E_\gamma))^2 + (\Delta E_e)^2 = k^2 E_\gamma + (\Delta E_e)^2\\
	&= a * E_\gamma + b \label{eq:grade}
\end{align}

aufgetragen (Abb. \ref{fig:fwhm_fit}) und eine Gradengleichung (Gl \eqref{eq:grade}) daran angepasst. Die Fitparameter sind dann $a = \SI{2.201(64)}{\electronvolt}$ und $b = \SI{8.95(19)e+05}{\electronvolt^2}$. Es ist dann die intrinsische Halbwertsbreite $\Delta E_d = \sqrt{b} = \SI{946(10)}{\electronvolt}$ und die Konstante $k = \sqrt{a} = \SI{1.482(22)}{\electronvolt}$.

\todo{kommentieren was das uns jetzt sagt...}
\jafps{fwhm_fit}{\todoP{write this}}{0.75}

